\section{Compact null families and descent}
\label{sg:descent}

We now complete Theorem~\ref{thm:signed}.  Retain $0<v<n$,
$r=v-1$, $d=n-1-r=n-v$, and the construction of
Proposition~\ref{prop:signed-construction}.  The formal lifts
first produce compact subvarieties outside the entire boundary,
of dimension $n-v$, on which $L$ is numerically trivial.
Their existence bounds the dimension of the nef reduction;
descent along that reduction then proves abundance.

The immediate geometric problem is to move a projective fiber off the
boundary while keeping it projective. A normal coordinate supplies the
direction of motion; parameters on the image of the boundary family keep
the fiber dimension fixed. The global root and Hodge constructions have
already supplied these functions formally. Their lifts concern the common
underlying topological map of the thickenings; a map of ringed spaces to
the base has not been assumed. The construction below obtains a deformation
directly from the lifted functions, without requiring one ambient scheme
neighborhood for all orders.

\begin{proposition}\label{prop:null-family}
Assume the transition surjectivity of
Proposition~\ref{prop:formal-lifting}.  There is a dominating
algebraic family of integral projective subvarieties of $X$
of dimension
\[
 d=n-1-r=n-v
\]
whose general members avoid $D$ and on which $L$ is
numerically trivial.
\end{proposition}

\begin{proof}
We construct a formal deformation of a general boundary fiber,
algebraize it over $\C[[s]]$, and show that the resulting
subvarieties belong to a family dominating $X$.

\smallskip\noindent\emph{Deforming a boundary fiber.}
Choose a separated affine etale chart $U\to\mathcal Y$
near a general point of a top-dimensional component of its
Stein image $\mathcal T_U$.  After shrinking, $C$ is trivial,
that component is smooth of pure dimension $r$, and $S_U$
is flat over it.  Since the image of $J$ in $P$ has dimension
less than $r$, we may also arrange that $S_U$ is disjoint
from $J_U$.  Choose a closed point $t$ in this open set.
After a further shrinking, regular parameters
$t_1,\ldots,t_r$ on $\mathcal T_U$ cut out only $t$.
The fiber
\[
 F=S_{U,t}
\]
is projective and has pure dimension $d$.  Along $F$, the
pullbacks of the parameters form a regular sequence.  The ideal
sheaf $\mathcal I_{J/S}$ is maximal Cohen--Macaulay by
Proposition~\ref{prop:signed-construction}(ii).  Since $S_U$ is
disjoint from $J_U$, the etale pullback of this ideal sheaf is
$\mathcal I_{J_U/S_U}=\OO_{S_U}$; hence $S_U$ is
Cohen--Macaulay.  The quotient $F$ by the regular sequence is
Cohen--Macaulay as well.  We will use this property to control
the dimensions of the components after algebraization.

Write $I=\OO_{\mathcal X}(-S)$.  By
Proposition~\ref{prop:formal-lifting}, the transitions between
the sheaves $g_*(I^j/I^{j+k})$ are surjective as the length
$k$ increases.  Their kernels are the coherent sheaves
$g_*\OO_{S_U}\otimes C^{-j-k}$.  Affineness of $U$ also gives
surjectivity on global sections.  We can therefore lift the
parameters compatibly through all thickenings, and lift the
chosen conormal frame to a compatible element $\widetilde y$
that formally generates $I$.

Let $Z_q$ be the closed subscheme of $qS_U$ cut out by the
$r$ lifted parameters.  Regard it as a scheme over
\[
 R_q=\C[s]/(s^q),\qquad s\longmapsto\widetilde y.
\]
The construction has the diagram
\[
\begin{tikzcd}[column sep=large]
F \arrow[r,hook] \arrow[d] & Z_q \arrow[r] \arrow[d] &
X\times\Spec R_q \arrow[d] \\
\Spec\C \arrow[r,hook] & \Spec R_q \arrow[r,equal] & \Spec R_q.
\end{tikzcd}
\]
The middle vertical map is defined by the lifted normal parameter.
The lifted base parameters cut out $Z_q$; they are not provided by a
preexisting morphism of the ambient thickening to the base.

The schemes $Z_q$ are compatible under reduction, have
closed fiber $F$, and are flat over $R_q$.  For flatness,
the powers of the Cartier generator first identify the
successive $s$-layers of $qS_U$ with $\OO_{S_U}$.  This
proves flatness before imposing the parameters.  The local
flatness criterion then preserves flatness when we quotient
by lifts of the regular sequence on the closed fiber.

\smallskip\noindent\emph{Algebraization and boundary avoidance.}
The map $Z_q\to X\times\Spec R_q$ is quasi-finite.  It is
also proper: its reduction is the map from the projective
scheme $F$, and nilpotent thickenings do not change
properness of a finite-type morphism.  Thus it is finite.
Put $R=\C[[s]]$.  Projective Grothendieck existence
algebraizes the compatible finite algebra sheaves on
$X\times\Spec R_q$ to a finite algebra on
$X_R=X\times\Spec R$.  Full faithfulness algebraizes
their multiplication and unit
\cite[Lemmas~30.24.1 and~30.24.3]{StacksExistence}.
The relative spectrum gives a finite morphism
\[
 Z\longrightarrow X_R
\]
with exactly the prescribed reductions.  In particular,
$Z$ is projective over $R$.  Formal flatness gives
flatness along the closed fiber; on the generic fiber
over $\C((s))$ it is automatic.  Hence $Z$ is flat over $R$.

The images of the pole and coefficient-zero boundary parts
miss $F$.  Their inverse images in $Z$ are closed and
proper over $R$, and must therefore be empty: any nonempty
closed subset of a proper $R$-scheme specializes to the
closed fiber.  Recall also the exceptional divisor $E$ of the
normalized blowup separating $S_+$ and $S_-$ in
Proposition~\ref{prop:signed-construction}.  Its image in $X$
lies in $\Supp G_+\cap\Supp G_-$, so it is already excluded
by avoidance of the pole part.

On the formal neighborhood from which the $Z_q$ were cut out,
the root construction therefore identifies the divisor of the
pulled-back section $s_0$ with $\ell S$.  Comparing its
root-line frame with the formal Cartier generator
$\widetilde y$ gives a frame in which this section is
$\widetilde y^\ell$.  Restricting to the $Z_q$ gives
compatible trivializations
\[
 N_0|_{\widehat Z}\simeq\OO_{\widehat Z},
 \qquad s_0|_{\widehat Z}=s^\ell.
\]
Full faithfulness algebraizes this line-bundle trivialization
and its section identity on $Z$; no algebraic map from $Z$
to the root stack is needed.  The generic fiber of $Z$ therefore misses
the positive part of $D$ as well, and hence avoids all of
$D$.

We next check the dimension assertion needed for the family.
At a closed point $z$ of $F$, flatness makes $s$ a
nonzerodivisor in $\OO_{Z,z}$, and its quotient is the
Cohen--Macaulay ring $\OO_{F,z}$ of dimension $d$.
Thus $\OO_{Z,z}$ is Cohen--Macaulay of dimension $d+1$.
These local rings are also equidimensional: a Cohen--Macaulay local ring
essentially of finite type over the excellent discrete
valuation ring $R$ has this property.
Every component of $Z$ meets $F$ by properness, and no
component is contained in $F$ by flatness.  The dimension
formula now gives dimension $d$ for every component of the
generic fiber.  Its finite image in $X_{\C((s))}$ therefore
has pure dimension $d$.  After a suitable finite extension of
$\C((s))$, choose a geometrically integral component of the
reduced base change of this image.

\smallskip\noindent\emph{A family dominating $X$.}
We have constructed compact subvarieties outside $D$.
To obtain a dominating family, choose a positive component
$D_i$ with image dimension $r$, together with a component
of $S$ covering it.  As the general chart and $t$ vary,
the images of points of $F$ contain a dense open subset
of $D_i$.  We will show that each point in this open subset
lies in the closed evaluation image of an algebraic family
of $d$-dimensional subvarieties disjoint from $D$.

Now consider all Hilbert schemes of $d$-dimensional
subvarieties of $X$.  The geometrically integral loci
parametrizing members disjoint from $D$ have a countable
stratification by integral parameter spaces with
irreducible universal families.  Denote the irreducible
closures of their evaluation images in $X$ by $W_\alpha$.
Fix a point $x$ in the dense open subset of $D_i$, and a
point of a boundary fiber $F$ mapping to it.  Since $Z$ has
no vertical component, this point lies in the closure of a
generic component.  That component has dimension $d$ by
the preceding argument.  Pass to a finite extension of the
discrete valuation field over which the reduced generic
components are geometrically integral, and let $R'$ be the
extended discrete valuation ring.  Choose a generic component
after this base change whose closure contains a point over
the chosen point of $F$.  Flatness still excludes vertical
components, so such a choice is possible.  Its reduced
finite image is geometrically integral.  Take the schematic
closure of that image in $X_{R'}$.
The structure sheaf of the closure has no uniformizer
torsion, so the closure is flat over the discrete valuation
ring and defines a morphism to the Hilbert scheme.
Its generic member lies in one of the strata above.
The corresponding $W_\alpha$ is closed and contains the
generic image, hence also its special-fiber support and
the point $x$.  Thus every point of the chosen open subset
of $D_i$ belongs to some $W_\alpha$.  The generic component
and the index $\alpha$ may depend on $x$.

An irreducible variety over the uncountable field $\C$
cannot have a dense open covered by countably many proper
closed subsets.  Some $W_\alpha$ therefore contains $D_i$.
By definition it also contains points outside $D$.
Since $D_i$ is a prime divisor in the integral variety
$X$, the only proper irreducible closed subset containing
it is $D_i$ itself.  Hence $W_\alpha=X$, and the family
dominates $X$.  On each member, avoidance of $D$ makes
the rational section $s_0$ regular and nowhere zero.
It trivializes $N_0$, proving that $L$ is numerically
trivial on that member.
\end{proof}

\subsection{Descent along the nef reduction}

The compact null subvarieties will bound the dimension of the
nef-reduction base.  To descend the signed representative to
that base, we first prove a statement for vertical divisors
using relative nefness and a fiber-dimension bound.

\begin{lemma}\label{sg:vertical-descent}
Let $f:Y\to B$ be a surjective projective morphism from a
normal integral variety to a smooth projective variety,
with geometrically connected integral generic fiber.
Suppose all fibers have dimension at most
$\dim Y-\dim B$.  Let $G_Y$ be a vertical $\Q$-Cartier
divisor that is relatively nef.  Then
\[
 G_Y=f^*G_B
\]
for a $\Q$-divisor $G_B$ on $B$.
\end{lemma}

\begin{proof}
When $B$ is a point, every vertical divisor is zero.
When the relative dimension is zero, the fiber bound
makes $f$ finite.  Its geometrically integral connected
generic fiber makes it birational, and normality of $B$
then makes it an isomorphism.  The assertion follows in
both cases.  We may therefore assume that the base
dimension and the relative dimension are positive.

The fiber bound forces every vertical prime divisor to
map onto a prime divisor of $B$.  For each such prime
$P$, write its full pullback as $f^*P=\sum_jm_jE_j$.
Let $b_j$ be the coefficient of $G_Y$ along $E_j$, with
zero coefficients included, and assign coefficient
$\min_j(b_j/m_j)$ to $P$ in $G_B$.  Only finitely many
of these coefficients are nonzero.  The residual divisor
\[
 R'=G_Y-f^*G_B
\]
is effective, vertical, $\Q$-Cartier, and relatively nef;
over each base prime, at least one component is absent
from its support.  It remains to prove that $R'=0$.

Suppose $R'\ne0$, and choose a base prime under its
support.  Intersect $B$ with general hyperplanes to
obtain a complete-intersection curve meeting that prime
at a general point; if $\dim B=1$, use $B$ itself.
Normal Bertini makes its inverse image in $Y$ normal.
This inverse image is integral: generic geometric
integrality supplies the dominating component, and
the fiber bound rules out further vertical components.
Intersect upstairs with $\dim Y-\dim B-1$ general
very ample hyperplanes.  The result is a normal integral
surface over the base curve, with geometrically connected
generic fiber.  Over the chosen general base point,
these cuts retain curves both in the residual support
and in a component missing from it.  The curves remain
distinct because the original components were distinct
at the generic point of the base prime.

Resolve this surface.  The pullback of the residual
divisor stays effective, vertical, and relatively nef.
It has zero intersection with the full fiber and
nonnegative intersection with every fiber component.
All those component intersections must therefore be
zero.  The fiber is connected, by generic connectedness
and Stein factorization over the normal base curve.
Its intersection matrix is negative semidefinite,
with kernel generated by the full fiber.  The residual
part over the chosen point is thus a multiple of the
full fiber.  Its missing component forces that multiple
to be zero, contradicting the retained component of
its support.
\end{proof}

\begin{proof}[Completion of the proof of Theorem~\ref{thm:signed}]
The cases $v=0$ and $v=n$ are covered by
Lemma~\ref{sg:numerical-boundary}.  For $0<v<n$,
Propositions~\ref{prop:signed-construction},
\ref{prop:formal-lifting}, and~\ref{prop:null-family}
give a dominating family of compact $L$-trivial
subvarieties of dimension $d=n-v$.

Apply the nef-reduction theorem to a Cartier multiple of
$L$ \cite[Theorem~2.1]{BauerEtAl2002}.  The resulting almost
holomorphic rational map has connected fibers.  The
divisor $L$ is numerically trivial on its compact general
fibers and has positive degree on every noncontracted
curve through a very general point of $X$.  Write $b$
for the base dimension.  A compact $L$-trivial $d$-fold
through such a point must be contracted: otherwise,
general ample slices through the point would produce a
noncontracted curve of $L$-degree zero.  Consequently
\[
 b\le n-d=v.
\]

To apply vertical descent, resolve the graph, flatten
the main component over a modification of the base,
resolve the base, and normalize the main transform.
This gives projective morphisms
\[
 \begin{tikzcd}
 Y \arrow[r,"\pi"] \arrow[d,"f"'] & X\\
 B &
 \end{tikzcd}
\]
with $\pi$ birational, $B$ smooth, $Y$ normal, and
geometrically connected integral generic fiber of $f$.
Every fiber has dimension at most $n-b$.  Indeed,
flatness and generic integrality keep the flat main
transform integral after the base modification, and
finite normalization preserves the bound on fiber
dimension.  This bound is the needed hypothesis;
normalization need not preserve flatness.

We claim that $\pi^*D$ has no horizontal component.
Restrict to a very general fiber of $f$.  There,
$\pi^*L$ is numerically trivial and $\pi^*(L-cD)$
restricts to a pseudo-effective class.  To justify the
second assertion, choose a sequence of effective approximants with
ample errors tending to zero. Shrink to a flat open of the base with
geometrically integral fibers. This is possible because the generic
fiber is geometrically integral. On this common open,
the locus over which a given support contains the whole fiber is
proper closed, by upper semicontinuity of fiber dimension.
Avoid the countably many such bad closed loci in the base.
Every approximant then restricts effectively to a very general fiber,
and passage to the limit gives pseudo-effectivity of the restriction.
The resulting class is $-c\pi^*D$.
A negative nonzero effective divisor cannot be
pseudo-effective on a projective variety, as its
intersection with an ample power shows.  The restriction
of $\pi^*D$ must therefore be zero.  This also excludes
$b=0$, which would force $D=0$ and contradict $v>0$.

It follows that $G_Y=\pi^*G$ is vertical.  It is also
relatively nef because $G_Y\sim_{\Q}\pi^*L$.
Lemma~\ref{sg:vertical-descent} now gives
\[
 \pi^*L\sim_{\Q}f^*G_B
\]
for a $\Q$-divisor $G_B$ on the smooth base.  Every
curve in $B$ is dominated by a curve in $Y$, where
the pullback has nonnegative degree; hence $G_B$ is
nef.  Numerical dimension is unchanged by these pullbacks;
the intersection formula gives
\[
 v=\nu(X,L)=\nu(Y,\pi^*L)=\nu(B,G_B)\le b.
\]
Together with $b\le v$, this yields $b=v=\nu(B,G_B)$.
Since $G_B$ is nef on the $b$-dimensional base, it follows
that $G_B^b>0$, so $G_B$ is big.  Pullback of
sections then gives $\kappa(X,L)\ge b=v$.  The
inequality $\kappa(X,L)\le\nu(X,L)$ for nef divisors
provides the reverse bound and proves abundance.
\end{proof}
